


\partwithintro[part:foundation]{The Foundation: Models for Brain Signals}{%
This part presents the foundation of the mind-world connection, focusing on the development of neuroscience-informed models for brain signals. We explore how these models can enhance interpretability, data efficiency, and performance on mind-related tasks by incorporating neural priors that reflect brain organization.}


\chapter{Learning High-Order Relationships of Brain Regions}

We begin our study of brain signals by using machine learning to discover candidate neural priors from data. In this chapter, we investigate high-order relationships among brain regions, which capture distributed patterns of brain activity beyond pairwise connectivity. We introduce \hybrid{}, a framework that learns these relationships from fMRI signals as weighted hyperedges. Through a multi-head drop-bottleneck, the model identifies relationships that are maximally informative about a phenotypic outcome while minimizing redundant information. The learned relationships improve phenotypic prediction and provide evidence for the value of high-order structure in modeling brain signals. This work establishes the discovery side of this part's foundation: machine learning can help identify patterns of neural organization that may inform model design. The next chapter examines the complementary direction, incorporating established spatial and anatomical priors into a neural encoder.

\label{chapter:hybrid}

\section{Introduction}
\label{sec:hybrid-intro}


Discovering relations among brain regions toward a specific phenotypic outcome from fMRI signals has been a crucial area in neuroimaging research. Reliable and informative relations help neuroscientists and clinical professionals to better understand brain functions, and thus improve clinical diagnosis and treatments \citep{kucian2008development,kucian2006impaired, li2015putting,li2015toward,satterthwaite2015linked,wang2016efficient}.
However, despite the clear multiplexity of the brain's involvement in cognition \citep{logue2014neural,barrasso2016neurobiological,knauff2010complex,reineberg2022context}, current imaging biomarker detection methods \citep{shen2017using,gao2019combining,li2021braingnn} focus only on the contributing roles of the pairwise connectivity edges. In contrast, most brain functions involve distributed patterns of interactions among multiple regions \citep{semedo2019cortical}. For instance, executive planning requires the appropriate communication of signals across many distinct cortical areas \citep{logue2014neural}. 
These high-order relationships, cannot always be decomposed into pairwise ones \cite{battiston2020networks, battiston2021physics, bick2023higher}, thus allowing them to capture information beyond the reach of pairwise ones \cite{do2020structural}. Consequently, relying solely on pairwise connectivity, without accounting for the brain's complex high-order structure, may result in inconsistent findings and low predictive performance across studies. 
Although recently there have been a few works \citep{zu2016identifying,xiao2019multi,li2022construction} working on discovering the high-order relationships of brain regions, they are unable to effectively extract meaningful patterns. This is because these methods usually first identify some candidates of high-order relationships and then perform feature selection. The candidates are obtained by enumerating only a small portion of all possible relations, or by clustering methods that are unrelated to the target. As a result, most informative high-order relations are likely excluded at the first stage. This inspires us to solve the problem in an end-to-end manner through a more expressive model and a more appropriate objective.

%1) the discovering process of high-order relationships is not guided by cognitive targets; 2) the weak expressiveness of traditional machine learning methods (e.g. lasso regression) hinders the accurate modeling of brain region interactions \citep{cao2022brain,richards2019deep}.

In this paper, we aim to identify high-order relationships that are informative towards a phenotypic outcome, such as a cognition score. However, unlike pairwise relations, the number of possible high-order relations is exponential. To identify the most informative ones from the exponential space, we propose our objective: \textit{maximally informative and minimally redundant} (MIMR). That is to say, we maximize the information contained in the high-order relationships towards a neurological outcome (\textit{informativeness}) while diminishing the participation of unrelated brain regions (\textit{redundant}).
% \todo{more citations in this paragraph to justify the principle would be awesome. We could even move some of the citations from previous paragraphs here if relevant.} 
Such a criterion, on the one hand, ensures the predictive performance of these high-order relationships; on the other hand, it endows the model with the capacity to identify more succinct and interpretable structures \citep{yu2020graph,miao2022interpretable,miao2022interpretable2,chen2023tempme}. A formal definition of the MIMR criterion could be found in Equation \ref{equ:ib2} from an information bottleneck point of view.

\begin{figure}
    \includegraphics[width=\linewidth]{figures/tasting400.png}
    % \vskip 0.1in
    \caption{We identify high-order relationships of brain regions, where hyperedge structures and weights possess strong relevance to cognition (maximal informativeness). Meanwhile, they contain the least irrelevant information.}
    \label{fig:tasting}
% \vskip -0.1in
    
\end{figure}

% This property, on the one hand, ensures the high-order relationships encapsulate sufficient information related to cognition, which prevents the omission of any of them that makes a significant contribution to cognition. On the other hand, it allows us to diminish irrelevant ones, which will make the results as succinct as possible.


% \htl{we (propose to) formulate the capture of high-order relationship in the brain as the construction of the hypergraph, where the nodes are regions and the hyperedges connect all the related regions.}
% \htl{To better...(require a motivation of why MRMI), a desired hypergraph should maximize the correlation between hyperedges and the predicted property while diminishing the correlation with unrelated inputs, which is also demonstrated by xxx.}

% \htl{However, most of the previous works on hypergraph construction fail to xxx since }

We formulate high-order relationships as weighted hyperedges in a hypergraph, where regions are treated as nodes. Unlike a traditional graph where edges connect only two nodes, a hypergraph allows edges, known as hyperedges, to connect any number of nodes. The hypergraph should be weighted, and the weights of hyperedges are considered as strengths of high-order relationships, which contain the information relevant to the outcome (Figure \ref{fig:tasting}). 
% Therefore, our goal is to construct a weighted hypergraph model of each brain, where the nodes are regions and hyperedges are set to be identified. These hyperedges should be maximally informative and minimally redundant: \textit{informativeness} is the mutual information between the hyperedges and the prediction target (e.g. cognition score).
% We maximize it because we expect these hyperedges to carry more information about the target; \textit{redundancy} is the mutual information between the hyperedges and the inputs. We minimize it because we expect succinct relations. Acquiring such maximally informative and minimally redundant hyperedges is much more challenging than pairwise edges in a normal graph due to the exponential search space since the number of nodes connected by a hyperedge can be arbitrary.\htl{better to separate it into two sentences "...due to...since"}

% In addition, \stodo{perhaps ad
% isk of overfitting in the model.

However, current methods for hypergraph construction, which are mostly based on neighbor distances and neighbor reconstruction \citep{wang2015visual,liu2017elastic,jin2019robust,huang2009video}, are unsuitable in our context for several reasons: 1) they are unable to learn MIMR hyperedges due to the absence of a tractable objective for learning such hyperedges.  2) they fall short of learning consistent structures across subjects, which contradicts the belief that the cognitive function mechanism of healthy human beings should be similar \citep{wang2023inference}. 3) the number of hyperedges is restricted to the number of nodes, which may lead to sub-optimal performance. Furthermore, although information bottleneck (IB) has been a prevailing solution to learn MIMR representations in deep learning \citep{kim2021drop,alemi2016deep,luo2019significance}, existing IB methods focus on extracting compressed representations of inputs instead of identifying underlying structures such as hypergraphs. Harnessing the IB framework for identifying hypergraphs necessitates both architectural innovations and theoretical derivations.

\paragraph{Proposed Work} In this paper, we propose \textbf{Hy}pergraph of \textbf{B}rain \textbf{R}egions via mult\textbf{i}-head \textbf{D}rop-bottleneck (\mname{}), a novel approach for identifying maximally informative yet minimally redundant high-order relationships of brain regions. The overall pipeline of \mname{} is depicted in Figure \ref{fig:pipeline}. \mname{} is equipped with a \textsc{Constructor} and a \textsc{Weighter}. The \textsc{Constructor} identifies the hyperedge structures of brain regions by learning sets of masks, and the \textsc{Weighter} computes a weight for each hyperedge. To advance the IB principle for hyperedge identification, we further propose \textit{multi-head drop-bottleneck} and derive its optimization objective.

\mname{} avoids searching in an exponential space through learning masks to identify hyperedges, which guarantees efficiency. Its feature-agnostic masking mechanism design ensures \mname{} to learn consistent structures across subjects. Moreover, the model is equipped with a number of parallel heads, each of which is dedicated to a hyperedge. Through this, \mname{} is able to identify any number of hyperedges, depending on how many heads it is equipped with. Additionally, the proposed \multiheadname{} theoretically guarantees the maximal informativeness and minimal redundancy of the identified hyperedges. 

% To alleviate the data scarcity issue, \mname{} is equipped with a special architecture, where some components are shared by all datasets, and others are specific to the condition.

We evaluate our methods on the open-source ABIDE dataset and the restricted ABCD dataset. We quantitatively evaluate our approach by a commonly used protocol for studying brain connections, CPM~\citep{shen2017using} (Appendix \ref{sec:app_cpm}), and show that our model outperforms the state-of-the-art deep learning models by an average of $11.2\%$ on a comprehensive benchmark. Our post-hoc analysis demonstrates that hyperedges of higher degrees are considered more significant, which indicates the significance of high-order relationships in human brains.
\section{Problem Definition \& Notations}
\label{sec:hybrid-prelims}
% \subsection{Problem Definition \& Notations}
\paragraph{Input} 
 Our dataset is a collection of human subject's features and their phenotypic outcomes, which is represented by the pair $(X, Y)$ for each subject. $X \in \mathbb{R}^{N \times d}$ represents the features for each subject, where $N$ is the number of brain regions and $d$ is the feature size. Consistent with previous works \citep{kan2022brain,li2021braingnn}, the features are Pearson correlations derived from fMRI time series. $Y \in \mathbb{R}$ denotes the phenotypic outcome, such as the intelligent quotient. Full implementation details and hyperparameter settings are provided in \cite{qiu2023learning}.
 
 %For how to obtain features of regions from raw data, see for more details.

\paragraph{Goal}
Based on the input $X$, \hybrid{} aims to learn a weighted hypergraph of the brain, where regions are nodes. To achieve this, \hybrid{} identifies a collection of hyperedges $H = (\boldsymbol{h}^1, \boldsymbol{h}^2, \cdots, \boldsymbol{h}^K)$, and assigns weights $\boldsymbol{w} = [w^1, w^2, \cdots, w^K]^T$ for all hyperedges. These hyperedges and their weights, which represent strengths of hyperedges, are expected to be the most informative towards $Y$ yet the least redundant. 

\paragraph{Representation of Hyperedges}
As mentioned before, we use $H$ to denote the collection of hyperedge structures and $\boldsymbol{h}^k$ to denote the $k$-th hyperedge. To associate hyperedges with node memberships, we use the following representation for a hyperedge:
\begin{equation}
    \boldsymbol{h}^k = \boldsymbol{m}^k \odot X \in \mathbb{R}^{N\times d},
\end{equation}
where $\boldsymbol{m}^k \in \{0,1\}^N$ is a mask vector and $\odot$ denotes broadcasting element-wise multiplication. In other words, each $\boldsymbol{h}^k$ is a randomly row-zeroed version of $X$.
\section{Related Work}
\paragraph{Hypergraph Construction}
 Existing hypergraph construction methods are mostly based on neighbor reconstruction and neighbor distances. For example, the $k$ nearest neighbor-based method \citep{huang2009video} connects a centroid node and its $k$ nearest neighbors in the feature space to form a hyperedge. \citet{wang2015visual,liu2017elastic,jin2019robust,xiao2019multi} further refine these neighbor connections through various regularization. However, the number of hyperedges of these methods is restricted to the number of nodes, and hyperedges obtained by these methods are inconsistent across instances. 
 Cluster-based approaches learn community structure in a graph \cite{bannadabhavi2023community,ying2018hierarchical}, which can also be used to form hyperedges.
 \citet{zhang2022deep,zhang2018dynamic} proposed to iteratively refine a noisy hypergraph, which is obtained by the aforementioned methods. Therefore, they share the same limitations as the aforementioned methods. In addition, these methods are unable to learn MIMR hyperedges due to the absence of a tractable objective. Other methods, such as attributed-based methods~\citep{huang2015learning, joslyn2019high}, are ill-adapted to our context since they require discrete labels. 
 %or a prior graph topology, which is usually very noisy.
 Different from these methods, we provide a way to learn a consistent hypergraph through a deep-learning model. Furthermore, thanks to the proposed \multiheadname, these hyperedges are theoretically ensured MIMR.

\paragraph{High-Order Relationships in fMRI}

Although there are some methods working on high-order relationships in fMRI, they have limitations and are inconsistent with our MIMR objective.
\citet{xiao2019multi,li2022construction} used the existing non-learning-based hypergraph construction methods, which may lead to noisy and inexpressive hypergraphs. \citet{zu2016identifying,santoro2023higher} enumerated all hyperedges with degrees lower than 3, which can only discover a tiny portion of all possible hyperedges in exponential space and is not scalable to a large degree. \citet{rosas2019quantifying} proposed O-information, which reflects the balance between redundancy and synergy. The O-information metric is utilized by \citet{varley2023multivariate} to study fMRI data. However, the objective of these methods is not consistent with ours: although both of us are quantifying the redundancy of high-order relations, our method is to learn those that are most informative toward a cognition score, while theirs is to depict the synergy and redundancy within a system.

\paragraph{Information Bottleneck}
Information bottleneck (IB) \citep{tishby2000information} is a technique in data compression. The key idea is to extract a summary of data, which contains the most relevant information to the objective. \citet{alemi2016deep} first employed an IB view of deep learning. After that, IB has been widely used in deep learning. The applications span areas such as computer vision \citep{luo2019significance,peng2018variational}, reinforcement learning \citep{goyal2019infobot,igl2019generalization}, natural language processing \citep{wang2020learning} and graph learning \citep{yu2020graph,yu2022improving, xu2021infogcl,wu2020graph}. Unlike these studies that use IB to extract a compressed representation or a select set of features, our approach focuses on identifying the underlying structures of the data.

\paragraph{Connectivity-based Phenotypic Prediction}
Recently, deep learning techniques have been increasingly employed in predicting phenotypic outcomes based on the connectivity of brain regions. Most works \citep{ahmedt2021graph,li2019graph,cui2022interpretable,kan2022fbnetgen,cui2022braingb,said2023neurograph} model the brain network as a graph, in which regions act as nodes and pairwise correlations form the edges. These methods predominantly utilize Graph Neural Networks (GNNs) to capture the connectivity information for predictions. In addition to GNNs, \citet{kan2022brain} proposed to use transformers with a specially designed readout module, leveraging multi-head attention mechanisms to capture pairwise connectivity. However, all of these methods heavily rely on pairwise connectivity and neglect more intricate higher-order relationships. This oversight, on the one hand, leads to sub-optimal prediction performances and, on the other hand, prevents domain experts from acquiring insightful neuroscience interpretations, given that brain functions often involves multiple regions.
\section{Method}
\label{sec:hybrid-approach}
\paragraph{Method Overview}

\hybrid{} consists of a \textsc{Constructor} $\mathcal{F}_c$, a \textsc{Weighter} $\mathcal{F}_w$, and a \textsc{LinearHead} $\mathcal{F}_l$. At a high level, the \textsc{Constructor} $\mathcal{F}_c$ is responsible for identifying hyperedges $H$ from the data to construct the hypergraph. After that, the \textsc{Weighter} $\mathcal{F}_w$ calculates a weight for each hyperedge. Finally, based on all the weights $\boldsymbol{w}$, the \textsc{LinearHead} $\mathcal{F}_l$ predicts the label $Y$.
An illustration of this pipeline is shown in Figure \ref{fig:pipeline}. The pipeline can be formulated as
\begin{equation}
% \vspace{-1em}
    \label{equ:pipeline}
    X \xrightarrow[\mathcal{F}_c]{} H \xrightarrow[\mathcal{F}_w]{} \boldsymbol{w}\xrightarrow[\mathcal{F}_l]{} Y.
\end{equation}

We will elaborate on the details of the architecture below.

\begin{figure*}
\vspace{0.8em}
    \centering
    \includegraphics[width=\linewidth]{figures/pipeline2.drawio.pdf}
    % \vskip 0.1in
    \caption{Overview of the \hybrid{} pipeline when the total number of hyperedges $K=3$. Hyperedge are in distinct colors for clarity.  The \textsc{Constructor} identifies hyperedges in the hypergraph, where regions are nodes. The \textsc{Weighter} computes a weight for each hyperedge. These weights, representing strengths of hyperedges, are expected to be informative in terms of our target $Y$. There are two separate phases after obtaining weights of hyperedges: 1) Training. The model's parameters are trained under the supervision of $Y$; 2) Evaluation. The output weights, as well as pairwise weights, are fed into the CPM \citep{shen2017using}.}
    \label{fig:pipeline}
    % \vskip -0.1in
\end{figure*}

\subsection{Learning the Hypergraph by Multi-head Masking}
\label{sec:learn-by-mask}

Each instance (human subject) is represented by $X = [X_1, X_2, \cdots, X_N]^T \in \mathbb{R}^{N\times d}$, where $X_i \in \mathbb{R}^d$ is a column vector representing the features of region $i$. These regions are nodes in the hypergraph we are going to construct. Hyperedges in the hypergraph can be beneficial in the learning process below because it is essential to model the relationships between more than two regions.

\paragraph{Hyperedges Construction}
In this paragraph, we elaborate on how the \textsc{Constructor} identifies the hyperedges, i.e. $H=\mathcal{F}_c(X)$.

Suppose the number of hyperedges is $K$, which is a predefined hyperparameter. We assign a head to each hyperedge. Each head is responsible for constructing a hyperedge by selecting nodes belonging to that hyperedge. 

Specifically, to construct the $k$-th hyperedge, the \textsc{Constructor}'s $k$-th head outputs a column vector $\boldsymbol{m}^k \in \{0,1\}^N$, where each element in the vector corresponds to a brain region,

\begin{equation}
\label{equ:shallow}
% \begin{aligned}
\boldsymbol{m}^k = [\mathbbm{1}(p^k_{\theta,1}), \mathbbm{1}(p^k_{\theta,2}), \cdots, \mathbbm{1}(p^k_{\theta,N})]^T \in \{0, 1\}^N,
% \end{aligned}
\end{equation}
where $p^k_{\theta, i} \in [0,1], i=1,2,\cdots, N$ are learnable probabilities. $\mathbbm{1}: [0,1] \mapsto \{0,1\} $ is an indicator function, which is defined as $\mathbbm{1}(x) = 1$ if $x > 0.5$ and $\mathbbm{1}(x) = 0$ if $x \leq 0.5$.  And $\boldsymbol{m}^k$ is a column vector corresponding to the $k$-th hyperedge. Note that since there is no gradient defined for the indicator operation, we employ the stop-gradient technique \citep{oord2017neural,bengio2013estimating} to approximate it.

\begin{figure}[t]
    \centering
    \includegraphics[width=0.63\linewidth]{figures/architecture.drawio.pdf}
    \caption{Architecture details of the \textsc{Constructor} and the \textsc{Weighter}  when the number of nodes $N=6$ and the number of hyperedges $K=3$. At a high level, the \textsc{Constructor} learns the hyperedge structure by masking nodes. The \textsc{Weighter} computes the weight of each hyperedge based on the hyperedge's member nodes and their features.}
    \label{fig:enter-label}
\end{figure}

In the vector $\boldsymbol{m}^k$, $0$ indicates nodes that are masked out, and $1$ indicates nodes that are not masked. Nodes that are not masked are considered to form a hyperedge together. We use $\boldsymbol{h}^k$ to represent the masked version of $X$
\begin{equation}
\begin{aligned}
    \boldsymbol{h}^k &= \boldsymbol{m}^k \odot X \\&= [\boldsymbol{m}^k_1 X_1, \boldsymbol{m}^k_2 X_2, \cdots, \boldsymbol{m}^k_N X_N] \in \mathbb{R}^{N\times d},   
\end{aligned}
\end{equation}
where $\odot$ is the broadcast element-wise multiplication. $\boldsymbol{m}^k_j$ is the $j$-th element of the vector $\boldsymbol{m}^k$.  

We obtain $K$ hyperedges for $K$ sets of masks. We use $H$ to denote the collection of all hyperedges.
\begin{equation}
    H = (\boldsymbol{h}^1, \boldsymbol{h}^2, \cdots, \boldsymbol{h}^K).
\end{equation} 

\paragraph{Hyperedge Weighting}
After obtaining the structure (i.e. member nodes) of each hyperedge,  the \textsc{Weighter} will calculate each hyperedge's weight, which is supposed to indicate the importance of that hyperedge, based on the member nodes and their features, i.e. $\boldsymbol{w} = \mathcal{F}_w(H)$.

These weights are obtained by a $\operatorname{Readout}$ module, which is composed of: 1) summing over all the non-masked nodes feature-wisely; 2) dim reduction operation.

\begin{equation}
\label{equ:readout}
    w^k = \operatorname{Readout}(\boldsymbol{h}^k) = \operatorname{DimReduction}({\boldsymbol{m}^k}^T\boldsymbol{h}^k) \in \mathbb{R},
\end{equation}

where $w^k$ is the weight of the $k$-th hyperedge and $\operatorname{DimReduction}$ is an MLP with ReLU activations, where the output dimension is $1$.
For all hyperedges, we obtain $K$ hyperedges in total, $\boldsymbol{w} = [w^1, w^2, \cdots, w^K]^T \in \mathbb{R}^K$. Finally, these weights will be fed into the final linear head to predict the label of the instance, 
\begin{equation}
\label{equ:linearhead}
    \hat{Y} = \mathcal{F}_l(\boldsymbol{w}) \in \mathbb{R}.
\end{equation}

In contrast to previous hypergraph construction methods \citep{jin2019robust,xiao2019multi}, which identify hyperedges by refining neighbors and simply aggregating node features, \hybrid{} makes these procedures learnable and thus is able to identify MIMR hyperedges in a data-driven way through expressive neural networks. The Full study that decides the number of hyperedges $K$ is provided in \cite{qiu2023learning}.

\paragraph{Computational Complexity}
The computational complexity of our model is $O(N^2K)$, which is just the same scale as that of MLPs even though we are addressing a more challenging task: identifying high-order relationships in an exponential space. Details of the complexity calculation can be found in \cite{qiu2023learning}


\subsection{Optimization Framework}
 Since there are no existing IB frameworks that can be applied in our context, we propose a new IB framework named \multiheadname{} to optimize \hybrid{}. To adopt an information bottleneck view of \hybrid{}, we consider $X$, $Y$ and $H$ are random variables in the Markovian chain $X \leftrightarrow Y \leftrightarrow H$. According to our MIMR objective, we optimize
\begin{equation}
\label{equ:ib2}
    \arg \max I(H; Y) - \beta I(H; X),
\end{equation}
where $I(\cdot ; \cdot)$ denotes the mutual information. $I(H; Y)$ corresponds to the informativeness and $I(H; X)$ corresponds to the redundancy. $\beta$ is a coefficient trading off informativeness and redundancy. Since optimizing the mutual information for high-dimensional continuous variables is intractable, we instead optimize the lower bound of Equation \ref{equ:ib2}. Specifically, for the first term (informativeness), it is easy to show
\begin{equation}
\label{equ:lowerbound}
\begin{aligned}
    I(H; Y) &= \mathbb{H}[Y] -  \mathbb{H}[{Y|H}] \\
    &=\mathbb{H}[Y] + \mathbb{E}_{p(Y, H)}[\log p(Y|H)] \\
    & = \mathbb{H}[Y] + \mathbb{E}_{p(Y, H)}[\log q_\phi(Y|H)] \\
    &\quad + \mathbb{E}_{p(H)}[\operatorname{KL}(p(Y|H)|q_\phi(Y|H))] \\
   & \geq \mathbb{H}[Y] + \mathbb{E}_{p(Y, H)}[\log q_\phi(Y|H)],
\end{aligned}
\end{equation}
where $\mathbb{H}[\cdot]$ is the entropy computation.
Since there is no learnable component in the entropy of $Y$, we only need to optimize the second term $\mathbb{E}_{p(Y, H)}[\log q_\phi(Y|H)]$. $q_\phi$ can be considered as a model that predicts $Y$ based on $H$, which essentially corresponds to $\mathcal{F}_l \circ \mathcal{F}_w$, where $\circ$ is the function composition. In practice, we set $q_\phi$ as a Gaussian model with variance $1$ as most probabilistic machine learning models do for continuous data modeling \citep{alemi2016deep,luo2019significance,peng2018variational,kingma2013auto}. 


For the second term (redundancy) in Equation \ref{equ:ib2}, we have

\begin{proposition}
    (Upper bound of $I(H;X)$ in \multiheadname{})
\begin{equation}
\label{equ:multi_drop}
    \begin{aligned}
        I(H; X) 
        \leq \sum_{k=1}^K
        I(\boldsymbol{h}^k; X)
        &\leq \sum_{k=1}^K \sum_{i=1}^N
        I(\boldsymbol{h}^k_i; X_i)\\
        &= \sum_{k=1}^K\sum_{i=1}^N \mathbb{H}[X_i] (1 - p^k_{\theta, i}),
    \end{aligned}
\end{equation}
\end{proposition}

where $\boldsymbol{h}^k_i$ and $X_i$  is the $i$-th row of $\boldsymbol{h}^k$ and $X$ respectively.  $p^k_{\theta,i}$ is the mask probability in Equation \ref{equ:shallow}. $\mathbb{H}$ is the entropy computation. The equality holds if and only if nodes are independent and hyperedges do not overlap. It's important to clarify that optimizing Equation \ref{equ:multi_drop} does not imply a penalty for overlaps. The second inequality is inspired by \cite{kim2021drop}. The proof of the proposition can be found in \citet{qiu2023learning}.


% \begin{equation}
%     \mathcal{F}_c(X) =  \mathbbm{1}(p_\theta) \odot X
% \end{equation}
% \begin{equation}
%     q_\phi (H) =\mathcal{F}_l \circ \mathcal{F}_w (H)
% \end{equation}
Therefore, instead of optimizing the intractable objective Equations \ref{equ:ib2}, we optimize its upper bound (i.e. loss function) according to Equations \ref{equ:lowerbound} and \ref{equ:multi_drop}.
\begin{equation}
\begin{aligned}
    \mathcal{L} &= \| Y - \mathcal{F}_l \circ \mathcal{F}_w \circ \mathcal{F}_c (X) \|_2^2 + \beta \sum_{k=1}^K\sum_{i=1}^N \mathbb{H}[X_i] (1 - p^k_{\theta,i}) \\&\geq - I(H, Y) + \beta I(H, X).
\end{aligned}
\end{equation}
In conclusion, The learnable components are the shallow embeddings in Equation \ref{equ:shallow}, the $\operatorname{DimReduction}$ MLP in Equation \ref{equ:readout} and the \textsc{LinearHead} $\mathcal{F}_l$ in Equation \ref{equ:linearhead}.
\section{Experiments}
\label{sec:hybrid-experiments}

In this section, we conduct experiments to validate the quality of the learned hyperedges in terms of the predictive performance towards the cognition phenotype outcome. Furthermore, we conduct ablation studies to validate the key components in our model. We also analyze our results both quantitatively and qualitatively.

\subsection{Predictive Performance of Hypereges}
% \begin{algorithm}
% \begin{algorithmic}
% \STATE    $d \gets \text{an random integer between } 2 \text{ and } 30. $

% \STATE $\mathcal{H} \gets \{\}$
% \FOR{$i = 1:d$}
% \STATE $h \gets \text{randomly sample } d \text{ nodes}$
% \STATE $\mathcal{H} \gets \mathcal{H} \cup \{h\}$
% \ENDFOR

% \FORALL{$h$ in $H$}
% \FORALL{node in $h$}

% \STATE $w \gets U(0, 1)$
% \STATE $x \gets U(0, 2w)$
% \STATE $y \gets y + $

% \ENDFOR
% \ENDFOR

% \end{algorithmic}

% \caption{The procedure to construct the synthetic dataset.}
% \label{alg:syn}
% \end{algorithm}


% \paragraph{Real-World Datasets}
\paragraph{Datasets}
\label{par:real_dataset}

We consider two fMRI datasets: 

1) \textit{Autism Brain Imaging Data Exchange (ABIDE)} \citep{craddock2013neuro} is an open-source dataset. This dataset involved resting-state fMRI of patients from 17 international sites, as well as the anatomical and phenotypic data. Regions are obtained by Craddock 200 atlas \citep{craddock2012whole}. We use the preprocessed version from the official website. For prediction targets, we choose three intelligence quotients: FIQ (full-scale intelligence), VIQ (verbal intelligence quotient), and PIQ (performance intelligence). 
%We aim to show our model's effectiveness on different cognition targets.

2) \textit{Adolescent Brain Cognitive Development (ABCD)} \citep{casey2018adolescent} is one of the largest public fMRI datasets. Access is limited and requires adherence to a rigorous data request procedure to acquire the data. The data is collected from $11,875$ children aged between 9 to 10 years old. The functional MRI (fMRI) data is collected from children when they were resting and when they performed three tasks (SST, EN-back, MID). We use the ABCD imaging data collected from the baseline (release 2.0) as well as the 2-year follow-up (release 3.0). In conclusion, we obtain $8$ sub-datasets (we refer to them as \textit{datasets} from now on) from $2$ timepoints under $4$ tasks. Regions are obtained by AAL3v1 atlas \citep{rolls2020automated}. For preprocess procedures of the ABCD dataset, please refer to \citet{qiu2023learning} for more details. For the prediction target, we consider fluid intelligence as our label. Fluid intelligence reflects the general mental ability and plays a fundamental role in various cognitive functions.
%We aim to show our model's effectiveness on different cognition targets
%Recent literature has seen success in predicting intelligence based on pairwise predictive modeling \citep{dubois2018resting,dubois2018distributed}, and we aim to improve the prediction accuracy of the current methods by involving high-order relationships on both resting-state and task-based data.

For region features, consistent with previous connectivity-based methods \citep{li2021braingnn,kan2022brain,ktena2018metric,said2023neurograph}, we use a region's Pearson correlation coefficients to all other regions as the region features.
% \paragraph{Features}
% \paragraph{Data preprocessing}
% \label{par:prepocessing}
% The raw fMRI data of an instance is represented in four dimensions (3 spatial dimensions + 1 temporal dimension), which can be imagined as a temporal sequence of 3D images. First, brain images are parceled into regions (or nodes) using the AAL3v1 atlas \citep{rolls2020automated}. Following previous works \citep{kan2022brain,li2021braingnn,thomas2022self}, each region's time series is obtained by averaging all voxels in that region. Consistent with previous connectivity-based methods \citep{li2021braingnn,kan2022brain,ktena2018metric,said2023neurograph}, for each region, we use its Pearson correlation coefficients to all regions as its features. We randomly split the data into train, validation, and test sets in a stratified fashion. The split ratio is 8:1:1.


\paragraph{Evaluation Metric}
To evaluate the quality of hyperedges obtained by \hybrid, we use CPM~\citep{shen2017using}, a standard model that could evaluate the relevance between the connectivity and the prediction target, due to its high impact in the community. 
In the original implementation of CPM, weights of pairwise edges are obtained by Pearson correlation between nodes. These weights, as pairwise connectivity, are fed into the CPM. CPM will output a metric that measures the overall correlation between edges and the prediction target, which can be considered as a measure of edge qualities. This process is formulated as 
\begin{equation}
\label{equ:metric0}
    r^\prime = \operatorname{CPM}(\boldsymbol{w}_p, Y),
\end{equation}
where $\boldsymbol{w}_p \in \mathbb{R}^{K_p}$ denotes the pairwise edge weights and $K_p$ is the total number of pairwise edges. $r^\prime$ is a metric that measures the quality of weights based on positive and negative correlations to the outcome. 

To evaluate the quality of the learned weights for our model, we replace the pairwise edge weights with the learned high-order weights $\boldsymbol{w}_h \in \mathbb{R}^{K_h}$, and thus adjust Equation \ref{equ:metric0} to 
\begin{equation}
\label{equ:metric}
    r = \operatorname{CPM}(\boldsymbol{w}_h, Y),
\end{equation}

% where $[\cdot \| \cdot]$ denotes the concatenation of two vectors. 
Comparing $r$ to $r^\prime$ reflects the quality of learned weights in terms of the prediction performance since it measures the overall correlation between weights and the prediction target. In our model, $\boldsymbol{w}_h = \boldsymbol{w}$, which is the learned hyperedge weights.

% The fMRI data is processed using BioImage Suite \citep{joshi2011unified} with the standard preprocessing procedures, such as slice time and motion correction, registration to the MNI template (see details in \cite{greene2018task} and \cite{horien2019individual}). We delete scans with more than 0.10 mm mean frame-to-frame displacement. 
% We use the ABCD imaging data collected from the baseline (release 2.0) as well as the 2-year follow-up (release 3.0). For each time point, we included children subjected to four conditions: 

% The fMRI data is processed using BioImage Suite \citep{joshi2011unified}. First, we performed motion correction and slice-time correction using SPM5; and via BioImage Suite, the data were registered to a standardized $3mm \times 3mm \times 3mm$ common space, where we generated masks representing white matter, gray matter, and cerebrospinal fluid (CSF) and computed the mean time courses for both white matter and CSF. We orthogonalized each gray matter time course with respect to the mean time courses of both white matter and CSF, and we orthogonalized each gray matter time course to the six motion-related signals via SPM5. We then applied a bandpass Butterworth filter with a frequency range of $0.02$Hz to $0.1$Hz to the orthogonalized time courses. We used a Gaussian kernel with a full-width at half-maximum (FWHM) of $6$mm to enhance spatial coherence and spatial smoothing. Lastly, we removed the linear trend from all signals in accordance with the methodology detailed in \cite{shenGroupwiseWholebrainParcellation2013}. We deleted scans with more than 0.10 mm mean frame-to-frame displacement. Additional details about the standard preprocessing procedures, such as slice time and motion correction, and registration to the MNI template can be found in \cite{greene2018task} and \cite{horien2019individual}. 

% For the prediction target, we consider fluid intelligence as our label. Fluid intelligence reflects the general mental ability and plays a fundamental role in various cognitive functions. Recent literature has seen success in predicting intelligence based on pairwise predictive modeling \citep{dubois2018resting,dubois2018distributed}, and we aim to improve the prediction accuracy of the current methods. 
\begin{table*}
\caption{$r$ values of our hyperedges compared to baselines on the ABCD dataset. Results are averaged over 10 runs. Deterministic methods do not have standard deviations.}
\label{tab:main}
% \vskip 0.1in
\resizebox{\linewidth}{!}{
\begin{tabular}{ccllllllll}
    \toprule
{Type} & {Model}& {SST 1} & {EN-back 1} & {MID 1}  & {Rest 1}& {SST 2} & {EN-back 2} & {MID 2} & {Rest 2} \\
         \midrule
Standard &pairwise& $0.113$ & $0.218$ &  $0.099$ & $0.164$  & $0.201$&$0.322$&$0.299$ & $0.289$ \\
\shline
\multirow{3}{*}{\shortstack{Hypergraph\\Construction}}
&$k$NN &  $0.115$ &  $0.268$ & $0.168$ &  $0.127$  & $0.257$&$0.266$&$0.238$ & $0.315$\\
&$l_1$ hypergraph  & $0.099$ &$0.223$  &  $0.125$ & $0.126$ &$0.145$&$0.295$&$0.242$& $0.259$\\
&$l_2$ hypergraph &$0.096_{\pm0.002}$  & $0.197_{\pm0.003}$ &  $0.118_{\pm0.003}$ & $0.157_{\pm0.016}$ &$0.203_{\pm0.005}$&$0.272_{\pm0.004}$&$0.289_{\pm0.011}$&$0.307_{\pm0.006}$\\
\shline
\multirow{4}{*}{\shortstack{Connectivity\\based\\Prediction}}
&BrainNetGNN &  $0.227_{\pm0.060}$ & $0.287_{\pm0.043}$ &$0.266_{\pm0.046}$ & $0.221_{\pm0.040}$ &$0.468_{\pm0.058}$&$0.480_{\pm0.068}$& $0.506_{\pm0.057}$&$0.453_{\pm0.028}$\\
&BrainGB & $0.190_{\pm0.073}$ & $0.214_{\pm0.051}$ &$0.265_{\pm0.048}$& $0.176_{\pm0.066}$ & $0.447_{\pm0.089}$ & $0.483_{\pm0.077}$ &$0.458_{\pm0.064}$ & $0.432_{\pm0.076}$\\
&BrainGNN & $0.262_{\pm0.030}$ & $0.235_{\pm0.032}$ &  $0.260_{\pm0.049}$ & $0.185_{\pm0.058}$ & $0.455_{\pm0.028}$ & $0.391_{\pm0.077}$&$0.445_{\pm0.078}$ &$0.368_{\pm0.041}$\\
&BrainNetTF &$\underline{0.327_{\pm 0.084}}$ & $\underline{0.338_{\pm 0.056}}$  & $\underline{0.370_{\pm 0.098}}$ & $\mathbf{0.334_{\pm 0.084}}$ & $\underline{0.633_{\pm0.178}}$& $\underline{0.631_{\pm0.142}}$ & $\underline{0.629_{\pm0.123}}$ & $\underline{0.588_{\pm0.138}}$ \\ 
\shline
Ours  & \hybrid & $\mathbf{0.361_{\pm0.058}}$  & $\mathbf{0.348_{\pm0.061}}$ & $\mathbf{0.386_{\pm0.060}}$ &  $\underline{0.223_{\pm0.056}}$ &$\mathbf{0.738_{\pm0.054}}$ & $\mathbf{0.714_{\pm0.037}}$&$\mathbf{0.816_{\pm0.053}}$ &$\mathbf{0.730_{\pm0.049}}$\\
         \bottomrule
    \end{tabular}}
% \vspace{-0.5em}
\end{table*}

 
\begin{table}[t]
 \caption{$r$ values of our hyperedges compared to baselines on the ABIDE dataset. Results are averaged over 10 runs. Deterministic methods do not have standard deviations.}
    \label{tab:abide}
% \vskip 0.1in
\centering
    % \resizebox{\linewidth}{!}{
    \begin{tabular}{cclll}
    \toprule
      {Type} & {Model} & {FIQ} & {VIQ} & {PIQ}  \\
      \midrule
        Standard & pairwise & $0.052$ & $0.124$ & $0.056$ \\
        \shline
        \multirow{3}{*}{\shortstack{Hypergraph\\Construction}}
        & $k$NN & $0.023$ & $0.093$ & $0.056$ \\
        & $l_1$ hypergraph & $0.043$ & $0.125$ & $0.061$ \\
        & $l_2$ hypergraph & $0.148_{\pm0.000}$ & $0.141_{\pm0.014}$ & $0.063_{\pm0.004}$ \\
        \shline
        \multirow{4}{*}{\shortstack{Connectivity\\based\\Prediction}}
        &BrainNetGNN & $\underline{0.162_{\pm0.042}}$ & $\underline{0.199_{\pm0.042}}$ & $\underline{0.223_{\pm0.025}}$ \\
        &BrainGB & $0.125_{\pm0.119}$ & $0.154_{\pm0.068}$ & $0.157_{\pm0.053}$ \\
        &BrainGNN & $0.105_{\pm0.041}$ & $0.176_{\pm0.049}$ & $0.159_{\pm0.051}$ \\
        &BrainNetTF & $0.132_{\pm0.111}$ & $0.176_{\pm0.053}$ & $0.180_{\pm0.054}$ \\ 
        \shline
        Ours  & \hybrid & $\mathbf{0.181_{\pm0.040}}$ & $\mathbf{0.204_{\pm0.031}}$ & $\mathbf{0.245_{\pm0.042}}$\\
     \bottomrule
    \end{tabular}
    % }
\end{table}

\begin{table*}[t]
\caption{Ablation studies on the masking mechanism. Results are averaged over 10 runs.}
% \vskip 0.1in
\label{tab:ablation_mask}
\resizebox{\linewidth}{!}{
\begin{tabular}{cccccccccccc}
    \toprule
{Model} & {SST 1} & {EN-back 1} & {MID 1} &{Rest 1}& {SST 2} & {EN-back 2} & {MID 2}  &{Rest 2}\\
         \midrule
 \hybrid  & $\mathbf{0.361_{\pm0.058}}$ &$\mathbf{0.348_{\pm0.061}}$ & $\mathbf{0.386_{\pm0.060}}$ & $\underline{0.223_{\pm0.056}}$ &$\mathbf{0.738_{\pm0.054}}$& $\mathbf{0.714_{\pm0.037}}$ &  $\mathbf{0.816_{\pm0.053}}$ & $\mathbf{0.730_{\pm0.049}}$\\
 \hybrid$_\mathrm{NoMask}$  & $0.297_{\pm0.035}$ & $0.274_{\pm0.057}$ & $\underline{0.323_{\pm0.059}}$& $0.221_{\pm0.034}$&$0.653_{\pm0.036}$&$0.599_{\pm0.059}$ &$0.757_{\pm0.021}$ &$0.543_{\pm0.038}$\\
 \hybrid$_\mathrm{RndMask}$  & $0.256_{\pm0.069}$ & $0.191_{\pm0.046}$ & $0.255_{\pm0.080}$ & $0.190_{\pm0.051}$ &$0.541_{\pm0.069}$& $0.514_{\pm0.038}$&$0.598_{\pm0.064}$ &$0.482_{\pm0.083}$\\
 \hybrid$_\mathrm{SoftMask}$  & $\underline{0.343_{\pm0.042}}$ & $\underline{0.314_{\pm0.040}}$ & $0.320_{\pm0.055}$& $\mathbf{0.245_{\pm0.061}}$ &$\underline{0.707_{\pm0.042}}$ &$\underline{0.662_{\pm0.058}}$&$\underline{0.796_{\pm0.031}}$ & $\underline{0.655_{\pm0.030}}$ \\
         \bottomrule
    \end{tabular}}
    % \vspace{-0.6em}
\end{table*}

\paragraph{Baselines}
We compare our method with $3$ classes of baselines: 1) \textit{standard} method, which is exactly the classical method that predicts outcomes based on pairwise edges \citep{shen2017using, DADI2019115, bhaa407}. The comparison with standard methods shows whether the high-order connectivity has its advantage over the classical pairwise one or not. 
2) \textit{hypergraph construction} methods. We consider $k$NN \citep{huang2009video}, $l_1$ hypergraph \citep{wang2015visual}, and $l_2$ hypergraph \citep{jin2019robust}.
3) \textit{connectivity-based  phenotypic prediction} methods, which are state-of-the-art predictive models based on brain connectivity. We consider BrainNetGNN \citep{mahmood2021deep}, BrainGNN \citep{li2021braingnn}, and BrainNetTF \citep{kan2022brain}. BrainGB \citep{cui2022braingb} is a study of different brain graph neural network designs and we include its best design as a baseline. Since none of these models are able to identify hyperedge structures of brain regions, we input their last layer embeddings (each entry as a weight) into the CPM model. Note that our weights $\boldsymbol{w}$ are also last layer embeddings in \hybrid{}.

% \subsection{Results}

% \paragraph{Quality of Hyperedges}
\paragraph{Results}
We report $r$ values by CPM in Table \ref{tab:abide} and Table \ref{tab:main}. As we can see, on the ABIDE dataset, \hybrid{} consistently outperforms all the baselines on different targets, with an average improvement of $8.8\%$ compared to the state-of-the-art model.
On the ABCD dataset, \hybrid{} outperforms the state-of-the-art predictive models on $7$ datasets of ABCD, with an average improvement of \avgimp{}. The results demonstrate our model is able to learn informative hyperedges towards different phenotypic outcomes from fMRI data of various brain states. Rest 1 is the only dataset that our model reaches the second-best.

Further, the comparison between our model and the pairwise baseline demonstrates the superiority of incorporating high-order relationships over relying solely on pairwise ones.

\paragraph{Runtime} \hybrid{} outperforms all other deep learning baselines in efficiency, with $87\%$ faster than the second-fastest one (BrainNetTF). Full runtime analysis is provided in \citet{qiu2023learning}.

\paragraph{Ablation Studies}
We conduct an ablation study on the effect of our masking mechanism. Specifically, we compare our model with $3$ variants: 1) \hybrid$_\mathrm{RndMask}$: Replace the learnable masks with random masks with the same sparsity, initialized at the beginning of training. 2) \hybrid$_\mathrm{NoMask}$: Do not mask at all, which means all nodes and their features are visible to each head. 3) \hybrid$_\mathrm{SoftMask}$: Remove the indicator function and use $p^k_{\theta, i}$ directly in Equation \ref{equ:shallow}. Ablation results are shown in Table \ref{tab:ablation_mask}. We find the original \hybrid{} and the \hybrid$_\mathrm{SoftMask}$ outperform all other variants, which demonstrates the effect of learnable masks. Moreover, the original \hybrid{} is better than its soft version \hybrid$_\mathrm{SoftMask}$, which demonstrates our sparse and succinct representations preserve better information than smooth ones.

\paragraph{Additional Experiments on Synthetic Dataset}
Since there are no ground-truth hyperedges in real-world datasets of learning informative hyperedges towards a specific outcome, we construct a synthetic dataset to verify if our model can recover the correct hyperedge structure under the MIMR objective. We use the precision, recall, and F1 score to measure the correctness of the learned hyperedges with respect to the ground truth. Although it is challenging to learn hyperedges when only supervised by the task label, we find that our model reaches high performances, with an average improvement of 28.3\% in terms of the F1 score, compared to the strongest baselines.


\subsection{Further Analysis}
% \subsubsection{Hyperedge Profiles}
In this subsection, we analyze the results of our model. We mainly use the ABCD dataset in the analysis since it is much larger than the ABIDE dataset.
\paragraph{Hyperedge Degree Distribution}
We plot the hyperedge degree distribution in Figure \ref{fig:degree_dist}. We find there are two distinct clusters in the figure. The first cluster is hyperedges with degree $\leq 5$. $1$-degree and $2$-degree hyperedges are special cases of our method: $1$-degree hyperedges are individual nodes, which imply the contribution of individual regions to the cognition. $2$-degree hyperedges reveal the importance of traditional pairwise connectivity. The other cluster concentrates around degree $25$, which implies the importance of relationships of multiple regions.

\begin{figure*}
% \vspace{-0.5em}
\begin{center}
    \begin{subfigure}[t]{0.32\textwidth}
        \includegraphics[width=\textwidth]{figures/degree_dist2.pdf}
        \caption{}
        \label{fig:degree_dist}
    \end{subfigure}
    \hspace{1em}
    \begin{subfigure}[t]{0.27\textwidth}
        \includegraphics[width=\textwidth]{figures/degree_sig.pdf}
        \caption{}
        \label{fig:degree_sig}
    \end{subfigure}
        \hspace{1em}
    \begin{subfigure}[t]{0.28\textwidth}
        \includegraphics[width=\textwidth]{figures/num_sig.pdf}
        \caption{}
        \label{fig:cmp_sig_pair}
    \end{subfigure}
    % \vskip 0.1in
    \caption{Hyperedge profiles. \textbf{(a)} Hyperedge degree distribution of learned hyperedges. \textbf{(b)} Correlation between hyperedge degree and significance. \textbf{(c)} Comparison between the number of hyperedges and pairwise edges under different significance thresholds. The total number of hyperedges is $32$. And the total number of pairwise edges is $26,896$.}
    \label{fig:hyperedge_profile}
    % \vskip -0.1in
    \end{center}
\end{figure*}

% To further understand the relationship between hyperedge degree and its significance, we next plot their correlation in Figure \ref{fig:degree_sig}. \jialin{can be removed}
% \begin{wrapfigure}{R}{0.49\textwidth}
%           \includegraphics[width=\linewidth]{figures/degree_sig.pdf}
%     \caption{Relationship between the significances and degrees of learned hyperedges.}
%     \label{fig:degree_sig}
% \end{wrapfigure}

\paragraph{Hyperedges with Higher Degree are More Significant}
Since CPM conducts a significance test on pairwise edges and hyperedges internally based on a linear regression model, we can obtain a P-value for each hyperedge from the significance test. We define the significance of a hyperedge as $1-P_v \in [0, 1]$ where $P_v$ is the P-value of that hyperedge.

The relationship between hyperedge degree and its significance is shown in Figure ~\ref{fig:degree_sig}. In this figure, we find a strong positive correlation between a hyperedge's degree and its significance, which indicates that interactions of multiple brain regions play more important roles in cognition than pairwise or individual ones. It is also worth mentioning that there is a turning point around degree $5$, which corresponds to the valley around $5$ in Figure \ref{fig:degree_dist}.

\paragraph{High-order relationships are Better than Pairwise Ones}
To compare the significance in cognition between pairwise edges and learned hyperedges, we plot the number of remaining edges under different thresholds in Figure \ref{fig:cmp_sig_pair}. We find out that the learned hyperedges are much more significant than pairwise ones. Also note that even if we set the threshold to an extremely strict value ($1\times 10^{-8}$), there are still $60\%$ hyperedges considered significant. This evidence shows that our high-order relationships are much more significant than the traditional pairwise connectivity, which implies relationships involving multiple brain regions could be much more essential in cognition.


\paragraph{Hyperedge Case Study} We visualize the most significant hyperedge of the EN-back task in Figure \ref{fig:hyperedge1}. We observe a coordinated interaction of numerous brain regions, each fulfilling specific roles. Notably, some of these regions serve multi-functional purposes:
\begin{figure*}
    \centering
    \includegraphics[width=\linewidth]{figures/hyperedge1.pdf}
    \caption{Visualization of the most significant hyperedge of the EN-back task.}
    \label{fig:hyperedge1}
\end{figure*}


\begin{itemize}
\item \textbf{Memory Processing} \textit{ParaHippocampal\_L}, \textit{Temporal Mid}: Essential for memory encoding and retrieval, these regions are integral to the EN-back task, facilitating the recall of previously viewed images.

\item \textbf{Emotional Processing} \textit{Amygdala\_R}: The amygdala is crucial for the processing of emotions, such as fear and pleasure. Since the EN-back task involves emotional stimuli, it is reasonable that the region is connected by the hyperedge.

\item \textbf{Visual Processing}: \textit{Calcarine\_R}, \textit{Lingual\_L}, \textit{Fusiform\_L}. These regions are responsible for visual perception and some of them are related to complex visual contents like symbols and human faces, which were presented during the fMRI task.

\item \textbf{Sensory} \textit{SupraMarginal\_L}: It is responsible for interpreting tactile sensors and perceiving limbs location. Its involvement is likely due to the requirement for participants to engage in specific physical actions, such as pressing buttons, during the task. \textit{Temporal Mid}: It functions in multi-modal sensory integration.

\item \textbf{Motor Control} \textit{Cerebellum}: It is primarily responsible for muscle control. \textit{Caudate}: It plays a crucial role in motor processes.  Its involvement is likely attributed to participants engaging in physical actions, like pressing buttons.

\item \textbf{Cognitive Control} \textit{ACC\_pre\_L}:  In the EN-back task, this region is likely crucial for maintaining focus, error detection and correction, conflict management in working memory, and modulating emotional responses to the task's demands.
\end{itemize}
\section{Conclusion}
In this work, we proposed \hybrid{} for identifying maximally informative yet minimally redundant (MIMR) high-order relationships of brain regions. To effectively optimize our model, we further proposed a novel information bottleneck framework and derived its theory. Our method outperforms state-of-the-art models. The result analysis shows the effectiveness of our model. We expect such advancements could benefit clinical studies, providing insights into neurological disorders, and offering improved diagnostic tools in neurology and other related fields.
\paragraph{Limitations}
\hybrid{} only considers static high-order relations. Given that fMRI tasks are dynamic, including temporal changes and interactions, it will be interesting to study the evolution of these high-order relationships.

Additionally, \hybrid{} does not offer a method for interpreting complex high-order relationships. This limitation is not specific to \hybrid{}, but is a common challenge in analyzing such relationships.



// TODO: FA-EEG

\import{SWA/}{main.tex}
